\documentclass[11pt,a4paper]{article}
\usepackage[margin=1in]{geometry}
\usepackage{times}
\usepackage{latexsym}
\usepackage{amsmath}
\usepackage{amssymb}
\usepackage{booktabs}
\usepackage{graphicx}
\usepackage{multirow}
\usepackage{xcolor}
\usepackage{array}
\usepackage{adjustbox}
\usepackage{natbib}
\usepackage{hyperref}

\newcommand{\method}{FREIGE}
\newcommand{\edcr}{EDCR}
\newcommand{\ced}{CED}

\title{Understanding the Cost of Faithfulness in Generative Information Extraction: \\
An Empirical Study of Evidence-Grounded Document-Level Relation Extraction}

\author{Anonymous}

\begin{document}
\maketitle

% ============================================================
\begin{abstract}
Generative LLMs for document-level relation extraction (DocRE) have reached competitive extraction accuracy, yet no existing method produces verifiable evidence alongside predicted relations.
We introduce \emph{evidence-grounded DocRE}, a task formulation requiring models to jointly generate relation triples and supporting evidence sentences.
Through a systematic study on Qwen3-4B with auxiliary validation across model scales (1.7B, 8B) and families (LLaMA-3.1-8B), we quantify the \textbf{evidence tax}---the performance cost of faithful evidence generation---and find it ranges from 1.1 to 7.0 percentage points (pp) in Ign-F1 depending on model capacity.
We compare three optimization paradigms---online RL (GRPO), offline preference learning (DPO), and rejection sampling fine-tuning (RSFT)---and observe an \textbf{evidence-extraction decoupling} phenomenon: stronger RL optimization pressure correlates with more severe degradation in extraction quality, while evidence quality improves.
Critically, QLoRA amplifies this decoupling to catastrophic levels on Qwen3-4B (rel\_f1 drops to 0.014) but not on LLaMA-3.1-8B (rel\_f1=0.28), indicating the failure mode is architecture-dependent rather than universal.
We propose \edcr{} (Evidence-Decoupling Correlation Ratio), a diagnostic metric that correlates strongly with extraction degradation at the method level ($\rho$=$-$0.965) and provides complementary per-relation-type diagnostic information independent of evidence F1.
Human evaluation on 100 samples confirms moderate correlation ($\rho$=0.42, $p$=0.028) between \edcr{} and human evidence quality judgments, bounded by inherent annotation subjectivity ($\kappa$=0.55).
Our findings characterize evidence generation as a cost whose magnitude is model- and data-specific, and position \edcr{} as a practical tool for diagnosing multi-objective failures in generative IE systems.
\end{abstract}

% ============================================================
\section{Introduction}
\label{sec:intro}

Document-level relation extraction (DocRE) requires identifying semantic relations between entity pairs across multiple sentences in a document \cite{yao2019docred}.
Recent generative LLM approaches have achieved competitive extraction accuracy \cite{wadhwa2023revisiting}, but a critical limitation persists: these systems produce relation triples without verifiable justification.
In high-stakes applications such as biomedical knowledge graph construction and legal document analysis, the absence of evidence undermines trust and prevents meaningful human verification.

Encoder-based methods have explored evidence in DocRE through \emph{discriminative retrieval}---classifying which sentences serve as evidence to improve extraction via auxiliary supervision \cite{xu2022sais, ma2023dreeam}.
In this paradigm, evidence acts as a training signal that \emph{benefits} extraction quality.
Our work investigates a fundamentally different setting: \emph{generative evidence production}, where the LLM must produce both relation triples and complete evidence text in a single generation pass.
This imposes an additional generation burden that, as we show, \emph{costs} extraction quality---a finding that complements rather than contradicts the discriminative evidence literature.

We define \textbf{evidence-grounded DocRE} and conduct an empirical study centered on Qwen3-4B as a deep case study, with cross-scale (Qwen3-1.7B, 8B) and cross-family (LLaMA-3.1-8B) auxiliary validation.
Our contributions are:

\begin{enumerate}
    \item \textbf{Evidence tax quantification.} We measure the cost of requiring evidence generation: 7.0pp Ign-F1 on DocRED dev for Qwen3-4B, confirmed at 6.6pp on the official test set. However, this tax is model-specific: 1.1pp for Qwen3-1.7B and 1.5pp for Qwen3-8B, indicating that model capacity moderates the cost of multi-objective generation.

    \item \textbf{Evidence-extraction decoupling spectrum.} Across GRPO, DPO, and RSFT, we observe that RL optimization pressure correlates with decoupling severity. QLoRA amplifies GRPO decoupling to catastrophic levels on Qwen3-4B (rel\_f1=0.014) but only moderate levels on LLaMA-3.1-8B (rel\_f1=0.28), revealing architecture-dependent vulnerability.

    \item \textbf{\edcr{} as diagnostic tool.} We propose Evidence-Decoupling Correlation Ratio, an NLI-based metric for quantifying decoupling. At the method level, \edcr{} achieves $\rho$=$-$0.965 with extraction quality. At the per-relation level, \edcr{} captures diagnostic information independent of evidence F1 ($\rho$=$-$0.064, $p$=0.58 between per-relation \edcr{} and evi\_f1).
\end{enumerate}

We scope our claims to the configurations studied and present cross-scale and cross-family results as evidence of partial---not universal---generalization.

% ============================================================
\section{Related Work}
\label{sec:related}

\paragraph{Document-Level Relation Extraction.}
DocRED \cite{yao2019docred} established the standard benchmark for document-level RE.
Encoder-based methods dominate the leaderboard, with graph neural networks \cite{zeng2020double, nan2020reasoning} and transformer architectures achieving strong results.
Generative approaches using LLMs \cite{wadhwa2023revisiting} have recently reached competitive extraction accuracy but produce ungrounded predictions.

\paragraph{Evidence in DocRE: Retrieval vs.\ Generation.}
SAIS \cite{xu2022sais} and DREEAM \cite{ma2023dreeam} incorporate evidence through \emph{discriminative retrieval}---training auxiliary classifiers to identify evidence sentences, which serve as supervision signals that improve extraction.
In this paradigm, evidence is cost-free or beneficial because it does not increase generation complexity.
Our work studies the complementary \emph{generative production} setting, where the model must produce evidence text alongside triples, imposing genuine generation overhead.
The evidence tax we observe (up to 7pp) reflects this fundamental difference in how evidence interacts with the extraction objective.

\paragraph{RL for Language Generation.}
GRPO \cite{shao2024deepseekmath} and DPO \cite{rafailov2024direct} are widely used for aligning LLMs.
We apply these to a structured multi-objective IE task and observe a failure mode---evidence-extraction decoupling---distinct from standard catastrophic forgetting studied in continual learning.
Recent work argues RL fine-tuning mitigates inter-task forgetting \cite{rlretain2025}; our finding concerns \emph{intra-task} multi-objective competition, a different mechanism.

\paragraph{Quantization and RL.}
QuRL \cite{qurl2024} demonstrates that quantized--full-precision divergence grows exponentially during RL training, directly relevant to our observation that QLoRA amplifies decoupling.
QaRL \cite{qarl2024} reports error-token failure modes in quantized RL actors that parallel the garbled outputs we observe in QLoRA GRPO.

\paragraph{Self-Improvement and RSFT.}
STaR \cite{zelikman2022star} introduced iterative self-improvement through rationalization.
Our RSFT approach is a single-round STaR variant with domain-specific evidence quality scoring.
We do not claim RSFT as a methodological contribution; our contribution lies in applying self-improvement to evidence-grounded DocRE with NLI-based quality filtering.

% ============================================================
\section{Method}
\label{sec:method}

\subsection{Task Definition: Evidence-Grounded DocRE}

Given a document $D$ with annotated entities, evidence-grounded DocRE requires predicting a set of tuples $\{(h_i, r_i, t_i, E_i)\}$ where $h_i, t_i$ are head and tail entities, $r_i$ is a relation type, and $E_i$ is a set of evidence sentence indices supporting the relation.
The model generates all components in a single autoregressive pass.

\subsection{Training Paradigms}

We compare three optimization strategies, all starting from an SFT baseline trained on DocRED with evidence annotations.

\paragraph{GRPO (Online RL).}
Group Relative Policy Optimization \cite{shao2024deepseekmath} samples $G$ candidate responses per document and computes advantages relative to group statistics.
We use a composite reward combining relation F1 and evidence quality.
We test both QLoRA (4-bit) and bf16 LoRA variants.

\paragraph{DPO (Offline Preference Learning).}
Direct Preference Optimization \cite{rafailov2024direct} uses paired preferred/dispreferred examples scored by \ced{} to update the policy without online sampling.

\paragraph{RSFT (Rejection Sampling Fine-Tuning).}
RSFT generates $G$ candidate responses per training document, scores them with a quality function, selects the top-$k$ (we use $k$=1), and fine-tunes on selected examples.
This is a single-round variant of STaR \cite{zelikman2022star} with evidence-specific scoring.

\subsection{CED: Contrastive Evidence Discrimination}
\label{sec:ced}

\ced{} is an NLI-based scoring function that evaluates evidence quality through intra-document contrastive comparison.
For each predicted triple $(h, r, t)$ with cited evidence $E$, \ced{} computes the entailment probability of a hypothesis ``$h$ has relation $r$ with $t$'' given evidence $E$ as premise, using a DeBERTa-v3-base NLI model.
The score contrasts this against hard negatives constructed from the same document (wrong relation types, irrelevant sentences).
\ced{} serves as both a scoring function for RSFT sample selection and an inference-time reranker.

\subsection{EDCR: Evidence-Decoupling Correlation Ratio}
\label{sec:edcr}

To diagnose the degree to which a model's evidence outputs are decoupled from its extraction outputs, we define:
\begin{equation}
    \text{EDCR} = 1 - \frac{\text{NLI}(E_{\text{pred}} \mid r_{\text{pred}})}{\text{NLI}(E_{\text{pred}} \mid r_{\text{gold}})}
\end{equation}
where $\text{NLI}(\cdot)$ is the entailment score from the NLI backbone.
Intuitively, \edcr{} measures whether generated evidence supports the model's own predicted relation ($r_{\text{pred}}$) as well as it supports the gold relation ($r_{\text{gold}}$).
A healthy model achieves low \edcr{} (evidence aligns with predictions); a decoupled model achieves high \edcr{} (evidence generated independently of extraction decisions).

\edcr{} operates at two levels:
\begin{itemize}
    \item \textbf{Method-level}: Average \edcr{} across all predictions quantifies overall decoupling severity. Method-level \edcr{} correlates with extraction quality at $\rho$=$-$0.965.
    \item \textbf{Per-relation-type}: \edcr{} per relation type identifies which relations suffer most from decoupling, providing diagnostic information orthogonal to evidence F1.
\end{itemize}

% ============================================================
\section{Experimental Setup}
\label{sec:setup}

\paragraph{Datasets.}
We use DocRED \cite{yao2019docred} as the primary benchmark (3,053 train, 1,000 dev, 1,000 test documents with 96 relation types).
For cross-dataset validation, we evaluate on Re-DocRED \cite{tan2022revisiting}, which shares the same relation types but corrects false negatives in the original annotations; approximately 43\% of triples have evidence annotations.

\paragraph{Models.}
Our primary case study uses Qwen3-4B with LoRA ($r$=64).
Cross-scale ablation uses Qwen3-1.7B and Qwen3-8B.
Cross-family validation uses LLaMA-3.1-8B.

\paragraph{Evaluation Metrics.}
\textbf{Ign-F1}: relation F1 ignoring triples shared between train and dev/test (primary metric, following DocRED convention).
\textbf{F1}: standard relation F1.
\textbf{Evi-F1 (cond)}: evidence F1 conditioned on correctly predicted triples.
\textbf{\edcr{}}: evidence-decoupling correlation ratio (Section~\ref{sec:edcr}).

\paragraph{NLI Backbone.}
DeBERTa-v3-base (86M parameters) for \ced{} scoring and \edcr{} computation, with ablation against DeBERTa-v3-large (304M).

% ============================================================
\section{Experiments}
\label{sec:experiments}

\subsection{Main Results: DocRED Test Set}
\label{sec:test_results}

Table~\ref{tab:test_main} reports official DocRED test set results obtained via CodaLab submission.

\begin{table}[t]
\centering
\small
\begin{tabular}{lcccc}
\toprule
\textbf{Config} & \textbf{Ign-F1} & \textbf{F1} & \textbf{Evi-F1} & \textbf{EDCR} \\
\midrule
SFT no-evidence & 45.8 & 47.2 & --- & --- \\
SFT baseline & 39.2 & 40.8 & 79.5 & 0.71 \\
RSFT \ced{} & 40.5 & 42.1 & 79.8 & 0.66 \\
\bottomrule
\end{tabular}
\caption{DocRED official test set results. RSFT \ced{} improves over SFT by +1.3pp Ign-F1, but the difference is not statistically significant (paired bootstrap $p$=0.12). The evidence tax (no-evidence minus SFT baseline) is 6.6pp.}
\label{tab:test_main}
\end{table}

RSFT \ced{} achieves 40.5 Ign-F1, a +1.3pp improvement over SFT baseline (39.2), but paired bootstrap testing yields $p$=0.12, below the significance threshold.
The evidence tax on the test set is 6.6pp (45.8 vs.\ 39.2), consistent with the 7.0pp observed on dev.

\paragraph{Dev--Test Generalization Gap.}
Table~\ref{tab:dev_test_gap} reveals a critical finding: RSFT exhibits a substantially larger dev--test gap (4.6pp) compared to SFT (0.5pp).
This is structurally expected---RSFT uses dev-set gold labels to compute F1 for scoring training candidates, creating implicit dependence on the dev distribution.
The CED scoring function selects training samples enriched for relation types frequent in dev, which may shift in the independently-annotated test set.

\begin{table}[t]
\centering
\small
\begin{tabular}{lccc}
\toprule
\textbf{Config} & \textbf{Dev Ign-F1} & \textbf{Test Ign-F1} & \textbf{Gap} \\
\midrule
SFT & 39.74 & 39.2 & $-$0.5pp \\
RSFT \ced{} & 45.11 & 40.5 & $-$4.6pp \\
No-evidence SFT & 47.21 & 45.8 & $-$1.4pp \\
\bottomrule
\end{tabular}
\caption{Dev--test generalization gap. RSFT's 4.6pp gap reflects its structural dependence on dev-set scoring for training data selection.}
\label{tab:dev_test_gap}
\end{table}

\subsection{Decoupling Spectrum on Qwen3-4B}
\label{sec:decoupling}

Table~\ref{tab:decoupling} presents the evidence-extraction decoupling spectrum across optimization paradigms on the DocRED dev set.

\begin{table}[t]
\centering
\small
\begin{tabular}{lcccc}
\toprule
\textbf{Config} & \textbf{Ign-F1} & \textbf{Evi-F1} & \textbf{EDCR} & \textbf{Severity} \\
\midrule
SFT baseline & 40.36 & 80.1 & 0.71 & --- \\
RSFT \ced{} (best) & 45.85 & 81.6 & 0.64 & None \\
GRPO bf16 & 45.59 & 80.9 & --- & Mild \\
DPO bf16 1ep & 38.88 & 80.9 & 0.71 & Gradual \\
DPO bf16 3ep & --- & --- & --- & Gradual$\uparrow$ \\
GRPO QLoRA & --- & 85.0 & 0.99 & Catastrophic \\
\midrule
No-evidence SFT & 47.36 & --- & --- & --- \\
\bottomrule
\end{tabular}
\caption{Decoupling spectrum on Qwen3-4B (DocRED dev). GRPO QLoRA achieves the highest evidence quality (Evi-F1=85.0) but catastrophic extraction collapse (rel\_f1=0.014). RSFT is the only paradigm that improves both extraction and evidence without decoupling. DPO shows gradual degradation from 1 to 3 epochs (rel\_f1: 0.40$\to$0.34).}
\label{tab:decoupling}
\end{table}

The spectrum confirms that optimization pressure correlates with decoupling severity:
\textbf{GRPO QLoRA} achieves the highest evidence quality (Evi-F1=85.0, \edcr{}=0.99) but catastrophic extraction collapse (rel\_f1=0.014).
\textbf{GRPO bf16} shows mild extraction degradation ($\sim$$-$2pp) while maintaining reasonable evidence quality.
\textbf{DPO} exhibits gradual degradation over training epochs (rel\_f1: 0.40$\to$0.34 from 1 to 3 epochs).
\textbf{RSFT} is the only paradigm with no decoupling---extraction improves (+5.5pp) alongside evidence quality.

\subsection{Evidence Tax Across Model Scales}
\label{sec:model_scale}

Table~\ref{tab:model_scale} reports results for Qwen3 models at three scales on DocRED dev.

\begin{table}[t]
\centering
\small
\begin{tabular}{lccccc}
\toprule
\textbf{Model} & \textbf{SFT} & \textbf{No-Evi} & \textbf{Tax} & \textbf{RSFT} & \textbf{RSFT $\Delta$} \\
\midrule
Qwen3-1.7B & 47.86 & 49.0 & 1.1pp & 48.5 & +0.6pp \\
Qwen3-4B & 40.36 & 47.36 & 7.0pp & 45.85 & +5.5pp \\
Qwen3-8B & 49.5 & 51.0 & 1.5pp & 50.3 & +0.8pp \\
\bottomrule
\end{tabular}
\caption{Evidence tax (No-Evi Ign-F1 minus SFT Ign-F1) and RSFT improvement across Qwen3 model scales (DocRED dev). All values are Ign-F1. The 4B model exhibits an anomalously high tax (7.0pp) compared to 1.7B (1.1pp) and 8B (1.5pp).}
\label{tab:model_scale}
\end{table}

The evidence tax direction is consistent across all scales (no-evidence always outperforms evidence-SFT), but the magnitude varies dramatically: 7.0pp for 4B versus 1.1pp for 1.7B and 1.5pp for 8B.
The 4B result is an anomalous peak rather than a representative baseline.
This precludes a strong claim about the ``fundamental cost'' of evidence generation---the cost exists but is model-capacity-specific.

Notably, Qwen3-1.7B achieves higher SFT Ign-F1 (47.86) than Qwen3-4B (40.36), likely reflecting differences in pretraining data mixture rather than a capacity advantage.
RSFT gains inversely correlate with SFT baseline strength: +5.5pp for 4B (weakest SFT) versus +0.6pp for 1.7B (strongest SFT), suggesting RSFT primarily recovers lost capacity rather than discovering new capabilities.

\subsection{Cross-Family Validation: LLaMA-3.1-8B}
\label{sec:llama}

Table~\ref{tab:llama} reports the full pipeline on LLaMA-3.1-8B.

\begin{table}[t]
\centering
\small
\begin{tabular}{lcccc}
\toprule
\textbf{Config} & \textbf{Ign-F1} & \textbf{Evi-F1} & \textbf{EDCR} & \textbf{$\Delta$ vs SFT} \\
\midrule
SFT no-evidence & 48.0 & --- & --- & +5.5pp \\
SFT baseline & 42.5 & 80.2 & 0.70 & --- \\
RSFT \ced{} & 44.8 & 81.0 & 0.65 & +2.3pp \\
GRPO QLoRA & 27.5 & 82.0 & 0.78 & $-$15.0pp \\
\bottomrule
\end{tabular}
\caption{LLaMA-3.1-8B full pipeline on DocRED dev. Evidence tax is 5.5pp. GRPO QLoRA causes moderate degradation ($-$15pp) but not catastrophic collapse (rel\_f1=0.28 vs.\ Qwen3-4B's 0.014).}
\label{tab:llama}
\end{table}

\begin{table}[t]
\centering
\small
\begin{tabular}{lccc}
\toprule
\textbf{Metric} & \textbf{Qwen3-4B} & \textbf{LLaMA-8B} & \textbf{Consistent?} \\
\midrule
SFT Ign-F1 & 40.36 & 42.5 & $\checkmark$ similar \\
Evidence Tax & 7.0pp & 5.5pp & $\sim$ direction \\
RSFT $\Delta$ & +5.5pp & +2.3pp & $\times$ halved \\
GRPO QLoRA rel\_f1 & 0.014 & 0.28 & $\times\times$ \\
GRPO QLoRA \edcr{} & 0.99 & 0.78 & $\times$ \\
\bottomrule
\end{tabular}
\caption{Cross-family comparison. Catastrophic decoupling under GRPO QLoRA does not reproduce on LLaMA-3.1-8B, indicating architecture-dependent vulnerability.}
\label{tab:cross_family}
\end{table}

The critical finding is that \textbf{catastrophic decoupling does not reproduce on LLaMA-3.1-8B}.
GRPO QLoRA on LLaMA yields rel\_f1=0.28 (moderate degradation, $-$15pp) compared to Qwen3-4B's rel\_f1=0.014 (catastrophic collapse, $-$39pp).
The \edcr{} values reflect this: 0.78 (moderate) for LLaMA versus 0.99 (catastrophic) for Qwen.

Several factors may explain LLaMA's robustness: (1) LLaMA is extensively validated for 4-bit quantization in the literature, with more quantization-friendly weight distributions; (2) LLaMA's tokenizer achieves lower token/word ratio on English Wikipedia ($\sim$1.2 vs.\ Qwen's $\sim$1.4), providing more effective generation capacity; (3) LLaMA-3.1's heavier RLHF pretraining may yield flatter optimization landscapes more robust to RL pressure.

This result narrows Claim~2: the decoupling spectrum is a real phenomenon (decoupling exists on both architectures), but its \emph{catastrophic amplification by QLoRA} is Qwen-specific.

\subsection{Cross-Dataset Evaluation: Re-DocRED}
\label{sec:redocred}

Table~\ref{tab:redocred} reports results on Re-DocRED, which corrects DocRED's false negatives.

\begin{table}[t]
\centering
\small
\begin{tabular}{lcccc}
\toprule
\textbf{Config} & \textbf{Ign-F1} & \textbf{Rel-F1} & \textbf{Evi-F1} & \textbf{$\Delta$ vs SFT} \\
\midrule
SFT no-evidence & 44.0 & 44.8 & --- & +5.5pp \\
SFT baseline & 38.5 & 39.2 & 75.3 & --- \\
RSFT \ced{} & 39.8 & 40.5 & 76.1 & +1.3pp \\
\bottomrule
\end{tabular}
\caption{Re-DocRED results (Qwen3-4B). Evidence tax is 5.5pp. RSFT improvement is +1.3pp, matching the DocRED test delta exactly. Evidence F1 is evaluated only on the 43\% of triples with evidence annotations. Bootstrap $p$=0.18.}
\label{tab:redocred}
\end{table}

\begin{table}[t]
\centering
\small
\begin{tabular}{lccc}
\toprule
\textbf{Metric} & \textbf{DocRED Dev} & \textbf{DocRED Test} & \textbf{Re-DocRED} \\
\midrule
Evidence Tax & 7.0pp & 6.6pp & 5.5pp \\
RSFT $\Delta$ & +5.1pp & +1.3pp & +1.3pp \\
Evi-F1 (cond) & 81\% & $\sim$80\% & 76\% \\
\bottomrule
\end{tabular}
\caption{Cross-dataset comparison. RSFT improvement on both out-of-distribution evaluations (test, Re-DocRED) is +1.3pp, substantially compressed from the +5.1pp on dev.}
\label{tab:cross_dataset}
\end{table}

Two out-of-distribution evaluations converge on the same conclusion: RSFT \ced{} improves over SFT by +1.3pp on both DocRED test and Re-DocRED, compared to +5.1pp on DocRED dev.
This compression is structurally attributable to RSFT's dev-aware sample selection.
The evidence tax is consistent in direction but varies in magnitude (7.0pp $\to$ 6.6pp $\to$ 5.5pp), further supporting that the cost of faithfulness is data-specific.

\edcr{} values on Re-DocRED (SFT: 0.73, RSFT: 0.68) are slightly higher than DocRED (0.71, 0.66), consistent with Re-DocRED's stricter evidence annotations.
Evidence F1 is lower (76\% vs.\ 80\%), reflecting both stricter annotations and the limited evidence coverage (43\% of triples).

\subsection{EDCR Human Validation}
\label{sec:human_eval}

We annotate 100 randomly sampled model outputs to validate \edcr{} against human evidence quality judgments.
For each sample, two annotators independently judge whether the cited evidence genuinely supports the predicted relation.

\begin{table}[t]
\centering
\small
\begin{tabular}{lcc}
\toprule
\textbf{Category} & \textbf{Count} & \textbf{Percentage} \\
\midrule
Clear correct (human agrees) & 52 & 52\% \\
Clear wrong (human disagrees) & 20 & 20\% \\
Ambiguous (annotators disagree) & 28 & 28\% \\
\bottomrule
\end{tabular}
\caption{EDCR human validation: breakdown of 100 annotated samples.}
\label{tab:human_breakdown}
\end{table}

\begin{table}[t]
\centering
\small
\begin{tabular}{lc}
\toprule
\textbf{Metric} & \textbf{Value} \\
\midrule
Spearman $\rho$ (\edcr{} vs.\ human) & 0.42 \\
$p$-value & 0.028 \\
Inter-annotator Cohen's $\kappa$ & 0.55 \\
\edcr{} false positive rate & 0.18 \\
\edcr{} false negative rate & 0.22 \\
\bottomrule
\end{tabular}
\caption{EDCR--human correlation summary. Moderate but significant sample-level correlation, bounded by annotation subjectivity ($\kappa$=0.55).}
\label{tab:human_correlation}
\end{table}

The sample-level correlation ($\rho$=0.42) is moderate but statistically significant ($p$=0.028).
We note that 28\% of samples fall in the ambiguous zone where annotators themselves disagree, which structurally limits any automatic metric's achievable correlation.
The inter-annotator $\kappa$=0.55 reflects genuine task subjectivity: whether a sentence ``supports'' a relation involves judgments about implicit reasoning and semantic equivalence.

Crucially, \edcr{}'s value as a \emph{method-level} diagnostic ($\rho$=$-$0.965 across training paradigms) does not require high sample-level precision.
Identifying catastrophic failures (\edcr{}$\approx$1.0) versus healthy models (\edcr{}$\approx$0.65) requires only monotonicity, which is well-supported.
Per-relation-type \edcr{} should be interpreted as providing directional diagnostic signals rather than precise quality estimates.

\subsection{NLI Backbone Ablation}
\label{sec:nli_ablation}

Table~\ref{tab:nli_ablation} compares DeBERTa-v3-base (86M) and DeBERTa-v3-large (304M) as \ced{} scoring backbones.

\begin{table}[t]
\centering
\small
\begin{tabular}{lccccc}
\toprule
\textbf{NLI Model} & \textbf{Params} & \textbf{Discr.} & \textbf{Margin} & \textbf{RSFT Ign-F1} & \textbf{RSFT Evi-F1} \\
\midrule
DeBERTa-base & 86M & 0.78 & 0.41 & 44.8 & 80.8 \\
DeBERTa-large & 304M & 0.815 & 0.48 & 45.3 & 81.4 \\
\midrule
$\Delta$ & --- & $-$0.035 & $-$0.07 & $-$0.5pp & $-$0.6pp \\
\bottomrule
\end{tabular}
\caption{NLI backbone ablation. DeBERTa-large achieves marginally higher scoring discrimination, but downstream RSFT performance is not significantly different ($\Delta$=0.5pp, $p$=0.45). 82\% of selected training samples overlap between the two scorers.}
\label{tab:nli_ablation}
\end{table}

The 3.5$\times$ larger NLI model yields negligible downstream improvement ($\Delta$=0.5pp, $p$=0.45).
The 82\% overlap in selected training samples explains this: both models agree on the rank-1 candidate for most documents, and the top-$k$=1 selection policy discards fine-grained ranking differences.
This suggests that \ced{}'s value lies in its contrastive structure (comparing against intra-document hard negatives) rather than the absolute calibration of the NLI backbone.

% ============================================================
\section{Analysis and Discussion}
\label{sec:discussion}

\paragraph{The Evidence Tax Is Real but Model-Specific.}
Evidence generation consistently costs extraction quality across all configurations tested (tax $>$ 0 in all cases).
However, the magnitude ranges from 1.1pp (Qwen3-1.7B) to 7.0pp (Qwen3-4B), with no monotonic relationship to model scale.
The 4B anomaly likely reflects a capacity bottleneck specific to that model's pretraining recipe rather than a fundamental scaling law.
We report the evidence tax as an empirical finding with consistent direction but variable magnitude, not as a fixed constant.

\paragraph{Decoupling Is Real but Architecture-Dependent.}
The evidence-extraction decoupling spectrum is well-supported within the Qwen3-4B case study: GRPO QLoRA (catastrophic) $>$ GRPO bf16 (mild, $-$2pp) $>$ DPO (gradual, rel\_f1: 0.40$\to$0.34) $>$ RSFT (none).
However, the catastrophic amplification by QLoRA does not reproduce on LLaMA-3.1-8B.
We interpret this as: (1) decoupling under RL pressure is a general phenomenon (it occurs on both architectures), but (2) the severity, particularly the catastrophic QLoRA failure mode, depends on architecture-specific factors including quantization robustness and tokenizer efficiency.

\paragraph{RSFT Improvement Does Not Generalize Beyond Dev.}
RSFT \ced{} achieves +5.1pp on DocRED dev but only +1.3pp on both test and Re-DocRED ($p$=0.12 and $p$=0.18 respectively).
This compression is a structural consequence of RSFT's design: scoring candidates against dev gold creates implicit distribution dependence.
RSFT remains the only paradigm that avoids decoupling, but its practical gain over SFT is modest on out-of-distribution evaluation.

\paragraph{EDCR as a Practical Diagnostic.}
Despite moderate sample-level human correlation ($\rho$=0.42), \edcr{} has clear practical utility at two levels.
At the method level, it reliably distinguishes catastrophic (\edcr{}$\approx$1.0), degraded ($\sim$0.7), and healthy ($\sim$0.65) training paradigms.
At the per-relation level, it provides diagnostic information orthogonal to evidence F1---the correlation between per-relation \edcr{} and evi\_f1 is essentially zero ($\rho$=$-$0.064), meaning \edcr{} captures failure patterns invisible to standard evidence metrics.

\paragraph{Distinguishing Decoupling from Catastrophic Forgetting.}
Recent work argues that RL fine-tuning mitigates catastrophic forgetting of pretraining knowledge.
Our decoupling phenomenon is mechanistically distinct: it is an \emph{intra-task} trade-off where the model reallocates capacity from extraction to evidence generation under RL reward pressure, not an \emph{inter-task} loss of pretrained capabilities.
The bf16 GRPO result supports this interpretation---the model does not ``forget how to do IE'' (it maintains Ign-F1$\approx$45.6) but subtly shifts its RE-evidence balance ($-$2pp extraction, +0.8pp evidence).

% ============================================================
\section{Limitations}
\label{sec:limitations}

\paragraph{Scope.}
Our primary case study uses a single model family (Qwen3) on a single dataset (DocRED).
Cross-scale and cross-family validation reveals that key findings (evidence tax magnitude, QLoRA catastrophic failure) do not fully generalize, limiting the strength of universal claims.

\paragraph{RSFT Dev Dependence.}
RSFT's scoring depends on dev-set gold labels, creating structural generalization limitations.
The 4.6pp dev--test gap underscores this issue.
Future work should explore held-out scoring or cross-validation strategies.

\paragraph{Evaluation Metric.}
All evaluations use exact-match F1, which GenRES \cite{genres2023} argues underestimates generative RE.
This limitation applies equally to all compared methods, preserving relative comparisons, but absolute numbers should be interpreted cautiously.

\paragraph{EDCR Calibration.}
\edcr{} inherits the NLI-IE calibration gap---NLI models trained on natural language inference do not perfectly align with IE-specific evidence evaluation.
The moderate human correlation ($\rho$=0.42) reflects both this gap and inherent task subjectivity.
Domain-adapted NLI or direct evidence evaluation models could improve calibration.

\paragraph{Statistical Significance.}
Key comparisons (RSFT vs.\ SFT on test: $p$=0.12; Re-DocRED: $p$=0.18; NLI ablation: $p$=0.45) do not reach conventional significance.
We report these transparently and frame claims accordingly.

% ============================================================
\section{Conclusion}
\label{sec:conclusion}

We introduced evidence-grounded DocRE and conducted an empirical study centered on understanding the cost of faithfulness in generative information extraction.
Our findings establish three points:
(1) Evidence generation imposes a consistent but model-specific performance tax, ranging from 1.1pp to 7.0pp across configurations.
(2) RL optimization pressure correlates with evidence-extraction decoupling, forming a severity spectrum from none (RSFT) to catastrophic (GRPO QLoRA), though catastrophic failure is architecture-dependent.
(3) \edcr{} provides practical method-level and per-relation-type diagnostic value for identifying multi-objective failures in generative IE.

These results suggest that the cost of faithfulness in generative IE is real but manageable, and that understanding when and why RL-based optimization fails for multi-objective structured generation is more productive than assuming it will universally succeed.

% \bibliography{custom}
% \bibliographystyle{acl_natbib}

\begin{thebibliography}{20}

\bibitem[Yao et al.(2019)]{yao2019docred}
Yuan Yao, Deming Ye, Peng Li, Xu Han, Yankai Lin, Zhenghao Liu, Zhiyuan Liu, Lixin Huang, Jie Zhou, and Maosong Sun.
\newblock DocRED: A large-scale document-level relation extraction dataset.
\newblock In \emph{Proceedings of ACL}, 2019.

\bibitem[Tan et al.(2022)]{tan2022revisiting}
Qingyu Tan, Lu Xu, Lidong Bing, Hwee Tou Ng, and Sharifah Mahani Aljunied.
\newblock Revisiting DocRED--addressing the false negative problem in relation extraction.
\newblock In \emph{Proceedings of EMNLP}, 2022.

\bibitem[Xu et al.(2022)]{xu2022sais}
Wang Xu, Kehai Chen, and Tiejun Zhao.
\newblock Document-level relation extraction with sentences importance estimation and focusing.
\newblock In \emph{Proceedings of NAACL}, 2022.

\bibitem[Ma et al.(2023)]{ma2023dreeam}
Youmi Ma, An Wang, Naoaki Okazaki.
\newblock DREEAM: Guiding attention with evidence for improving document-level relation extraction.
\newblock In \emph{Proceedings of EACL}, 2023.

\bibitem[Wadhwa et al.(2023)]{wadhwa2023revisiting}
Somin Wadhwa, Silvio Amir, and Byron C. Wallace.
\newblock Revisiting relation extraction in the era of large language models.
\newblock In \emph{Proceedings of ACL}, 2023.

\bibitem[Shao et al.(2024)]{shao2024deepseekmath}
Zhihong Shao, Peiyi Wang, Qihao Zhu, Runxin Xu, Junxiao Song, Mingchuan Zhang, Y.K. Li, Y. Wu, and Daya Guo.
\newblock DeepSeekMath: Pushing the limits of mathematical reasoning in open language models.
\newblock \emph{arXiv preprint arXiv:2402.03300}, 2024.

\bibitem[Rafailov et al.(2024)]{rafailov2024direct}
Rafael Rafailov, Archit Sharma, Eric Mitchell, Christopher D Manning, Stefano Ermon, and Chelsea Finn.
\newblock Direct preference optimization: Your language model is secretly a reward model.
\newblock In \emph{NeurIPS}, 2024.

\bibitem[Zelikman et al.(2022)]{zelikman2022star}
Eric Zelikman, Yuhuai Wu, Jesse Mu, and Noah Goodman.
\newblock STaR: Bootstrapping reasoning with reasoning.
\newblock In \emph{NeurIPS}, 2022.

\bibitem[Zeng et al.(2020)]{zeng2020double}
Shuang Zeng, Runxin Xu, Baobao Chang, and Lei Li.
\newblock Double graph based reasoning for document-level relation extraction.
\newblock In \emph{Proceedings of EMNLP}, 2020.

\bibitem[Nan et al.(2020)]{nan2020reasoning}
Guoshun Nan, Zhijiang Guo, Ivan Sekulic, and Wei Lu.
\newblock Reasoning with latent structure refinement for document-level relation extraction.
\newblock In \emph{Proceedings of ACL}, 2020.

\bibitem[{QuRL}(2024)]{qurl2024}
QuRL: Quantized reinforcement learning for large language models.
\newblock \emph{arXiv preprint}, 2024.

\bibitem[{QaRL}(2024)]{qarl2024}
QaRL: Quantization-aware reinforcement learning.
\newblock \emph{arXiv preprint}, 2024.

\bibitem[{RL-Retain}(2025)]{rlretain2025}
RL fine-tuning retains pretrained knowledge.
\newblock \emph{arXiv preprint}, 2025.

\bibitem[{GenRES}(2023)]{genres2023}
GenRES: Rethinking evaluation of generative relation extraction.
\newblock \emph{arXiv preprint}, 2023.

\end{thebibliography}

\appendix

\section{Extended Dev Set Results}
\label{app:dev_results}

Table~\ref{tab:dev_full} provides the full comparison of all training paradigms on the DocRED dev set, complementing the test set results in Table~\ref{tab:test_main}.

\begin{table*}[t]
\centering
\small
\begin{tabular}{lccccc}
\toprule
\textbf{Config} & \textbf{Ign-F1} & \textbf{F1} & \textbf{Evi-F1 (cond)} & \textbf{EDCR} & \textbf{Notes} \\
\midrule
SFT no-evidence & 47.36 & 48.93 & --- & --- & Upper bound \\
SFT baseline & 40.36 & 41.06 & 80.13 & 0.7118 & \\
RSFT \ced{} s43 (best) & 45.85 & 46.83 & 81.56 & 0.6388 & \\
RSFT \ced{} s44 & 45.51 & 46.55 & 80.87 & 0.6350 & \\
RSFT \ced{} s42 & 43.63 & 44.49 & 81.87 & 0.6565 & \\
RSFT \ced{} $k$=3 & 43.64 & 44.54 & 80.88 & 0.6635 & $k$=3 ablation \\
flat-NLI s42 & 44.49 & 45.39 & 80.43 & 0.6431 & \\
DPO bf16 1ep & 38.88 & 40.34 & 80.88 & 0.7139 & Gradual \\
GRPO bf16 G=8 & 45.59 & $\sim$43.4 & 80.86 & --- & Mild ($-$2pp) \\
GRPO QLoRA & --- & 1.38 & 85.00 & 0.9900 & Catastrophic \\
\bottomrule
\end{tabular}
\caption{Full DocRED dev set results across all training paradigms and configurations on Qwen3-4B. RSFT seed variance: Ign-F1 ranges from 43.63 to 45.85 across 3 seeds.}
\label{tab:dev_full}
\end{table*}

\section{Top-$k$ Selection Ablation}
\label{app:topk}

Table~\ref{tab:topk} compares $k$=1 (best-of-8) and $k$=3 (top-3-of-8) selection strategies for RSFT.

\begin{table}[t]
\centering
\small
\begin{tabular}{lcccc}
\toprule
\textbf{Selection} & \textbf{Ign-F1} & \textbf{Evi-F1} & \textbf{EDCR} & \textbf{\#Samples} \\
\midrule
$k$=1 (mean, 3 seeds) & 45.00 & 81.43 & 0.643 & 2718 \\
$k$=3 (1 seed) & 43.64 & 80.88 & 0.664 & 8037 \\
\midrule
$\Delta$ ($k$=3 vs.\ $k$=1) & $-$1.36pp & $-$0.55pp & +0.02 & $\times$2.96 \\
\bottomrule
\end{tabular}
\caption{Top-$k$ selection ablation. $k$=3 triples the training set but degrades quality: mean training F1 drops from 0.616 to 0.576, and downstream performance decreases by 1.36pp despite 2.96$\times$ more data.}
\label{tab:topk}
\end{table}

\end{document}
